\documentclass[11pt,a4paper]{article}

\usepackage[T1]{fontenc}
\usepackage[utf8]{inputenc}
\usepackage[margin=24mm,headheight=15pt]{geometry}
\usepackage{amsmath,amssymb,mathtools}
\usepackage{newpxtext,newpxmath}
\usepackage{microtype}
\usepackage{xcolor}
\usepackage{needspace}
\usepackage{titlesec}
\usepackage{fancyhdr}
\usepackage{hyperref}
\usepackage{bookmark}

\definecolor{ink}{HTML}{243E49}
\definecolor{accent}{HTML}{176D74}
\definecolor{reference}{HTML}{844B3C}
\definecolor{muted}{HTML}{59686D}
\hypersetup{
  colorlinks=true,linkcolor=accent,urlcolor=reference,
  pdftitle={Mathematical Gems from OT4ML: 75 Possible Blog Posts},
  pdfauthor={Gabriel Peyre},
  pdfsubject={An editorial catalogue of mathematical blog-post ideas from OT4ML},
  bookmarksnumbered=false
}
\urlstyle{same}
\setlength{\parindent}{0pt}
\setlength{\parskip}{3pt}
\setlength{\emergencystretch}{1.5em}
\setlength{\abovedisplayskip}{7pt plus 2pt minus 2pt}
\setlength{\belowdisplayskip}{7pt plus 2pt minus 2pt}
\linespread{1.035}
\raggedbottom
\widowpenalty=10000
\clubpenalty=10000

\pagestyle{fancy}
\fancyhf{}
\fancyhead[L]{\small\color{muted}Mathematical Gems from OT4ML}
\fancyhead[R]{\small\color{muted}Blog-post catalogue}
\fancyfoot[C]{\small\color{muted}\thepage}
\renewcommand{\headrulewidth}{0.25pt}
\titleformat{\section}{\Large\bfseries\color{ink}}{}{0pt}{}
\titlespacing*{\section}{0pt}{18pt}{8pt}

\newcounter{post}
\newcommand{\post}[1]{%
  \par\addvspace{10pt}%
  \Needspace{7\baselineskip}%
  \begin{minipage}{\linewidth}%
  \refstepcounter{post}%
  \pdfbookmark[1]{\thepost. #1}{post.\thepost}%
  {\large\bfseries\color{accent}\thepost.\enspace #1\par}%
  \nobreak\vspace{2pt}\nobreak
}
\newcommand{\bookref}[3]{%
  \href{https://www.gpeyre.com/ot4ml/OT4ML/OT4ML.pdf\#nameddest=#3}{#1, p.~#2}%
}
\newcommand{\source}[1]{%
  \par\nobreak{\small\color{muted}\textit{Book pointer:} #1.\par}%
  \end{minipage}\par
}
\newcommand{\illustration}[1]{\par\textit{Possible illustration.} #1}
\newcommand{\Wass}{\mathcal W}
\newcommand{\SW}{\mathop{\mathrm{SW}}}
\newcommand{\GW}{\mathop{\mathrm{GW}}}
\newcommand{\MinSW}{\mathop{\mathrm{Min\text{-}SW}}}
\newcommand{\KL}{\mathop{\mathrm{KL}}}
\newcommand{\Law}{\operatorname{Law}}
\newcommand{\Id}{\operatorname{Id}}
\newcommand{\tr}{\operatorname{tr}}
\newcommand{\E}{\mathbb E}
\newcommand{\R}{\mathbb R}
\newcommand{\N}{\mathcal N}
\newcommand{\norm}[1]{\lVert #1\rVert}

\begin{document}
\thispagestyle{empty}
{\small\sffamily\color{accent}EDITORIAL NOTEBOOK\hfill 15 October 2026\par}
\vspace{13pt}
{\Huge\bfseries\color{ink}Mathematical Gems\par}
\vspace{3pt}
{\Huge\bfseries\color{ink}from OT4ML\par}
\vspace{9pt}
{\Large\color{muted}75 possible blog posts\par}
\vspace{15pt}

This catalogue collects all 75 ideas proposed after screening
\href{https://www.gpeyre.com/ot4ml/}{\emph{Optimal Transport for Machine Learners}},
in the spirit of the mathematical posts on
\href{https://www.gpeyre.com/blog/}{Gabriel Peyr\'e's blog}.
The editorial filter is simple: \emph{one surprising statement, a proof with an
identifiable mechanism, and a picture that reveals something the formula conceals}.

These are post proposals, not completed articles or claims of novelty.
Some are classical results with an underexposed interpretation; others are
counterexamples or synthesis pieces that deserve a more substantial mathematical
note. The hypotheses and attributions in the linked book passages should be
retained when a proposal becomes a post.

\section*{Where to Start}
\pdfbookmark[0]{Where to Start}{start}
\textbf{A first sequence: 1, 6, 10, 12 and 61.}
These combine concise proofs, distinctive mechanisms and straightforward visual
or interactive illustrations: inversion penalties, lattice effects in the CLT,
maximal correlation in Sinkhorn, failure of Hellinger displacement convexity, and
the logarithmic mean in Markov diffusion.

\textbf{Deeper mathematical notes: 2, 3 and 4.}
Scalar moment closure, non-Gaussian GW optimizers and the topology of a
measure-level Min-SW extension are particularly promising. Their exact scope and
literature positioning deserve extra care before publication.

\textbf{Avoiding repetition.}
The list favors new angles over direct repetitions of existing posts on Gaussian
preservation, Gaussian KL barycenters, Perron--Frobenius/Sinkhorn, moment measures,
complete positivity, inverse gap losses and Muon. Related proposals are intended
as sequels, contrasts or new connections.

\textbf{How to use the pointers.}
Each entry has a printable theorem or section locator and a clickable link into
the online book PDF. Numbering and printed page numbers refer to the current
manuscript snapshot; the source is standalone and does not depend on the book's
LaTeX files. PDF bookmarks give direct access to all 75 proposals.

\clearpage
\section*{Priority Candidates}
\pdfbookmark[0]{Priority Candidates: 1--12}{priority}

\post{The Monge gap is a sum of penalties for inversions}
On a uniform grid $x_i=x_1+(i-1)\Delta$, let $t_i=T(x_i)$ and
$\alpha_n=n^{-1}\sum_i\delta_{x_i}$. For quadratic cost,
\[
  \mathcal M_{\alpha_n}^{|\cdot|^2}(T)
  =\frac{2\Delta}{n}\sum_{1\leq i<j\leq n}(t_i-t_j)_+.
\]
A comparison with optimal transport becomes a convex penalty for being out of
order. Bubble sort explains the identity. The post can distinguish rearrangement
(which preserves the output law), isotonic regression (which pools values), and
Monge-gap-penalized regression.
\illustration{Reuse the kernel-regression regularization paths, with inversions highlighted.}
\source{\bookref{Proposition~3.44}{62}{thm.3.4.44}}

\post{Which statistics have an autonomous Wasserstein evolution?}
For $f(\alpha)=g(\int\varphi\,d\alpha)$, the evolution of the scalar moment
closes for every smooth $g$ and every admissible initial law precisely when
\[
  |\nabla\varphi|^2=a+b\varphi.
\]
The necessity follows by comparing two-point mixtures with equal moments, which
forces a Jensen functional equation. The condition connects moment closure to
the eikonal equation. Signed distance to a curved hypersurface supplies a
non-polynomial, locally valid example, as long as the support remains in a
smooth tubular neighborhood.
\source{\bookref{Proposition~16.30}{426}{thm.16.5.30}}

\post{Gaussian marginals need not have a Gaussian GW optimizer}
For the distortion $(|x-x'|^2-|y-y'|^2)^2$, couple a one-dimensional standard
Gaussian to a sufficiently high-dimensional standard Gaussian. A coupling that
correlates the radii can beat every jointly Gaussian coupling. Its advantage
comes from a mixed fourth moment that covariance optimization cannot detect.
This sharply separates Gaussian-constrained GW from unrestricted GW; the
dimensional and dissimilarity assumptions must remain explicit.
\illustration{Plot the Gaussian optimum and the radial competitor against dimension.}
\source{\bookref{Proposition~13.34}{300}{thm.13.3.34}; discuss the Gaussian-optimality
conjecture of Delon, Desolneux and Salmona cited there}

\post{An exact Gaussian formula can coexist with empirical inconsistency}
For the representation-invariant measure-level Min-SW extension used in the book,
$\MinSW_2=\Wass_2$ on nondegenerate Gaussians. Yet in dimension $d\geq2$,
empirical measures can converge to the standard Gaussian in $\Wass_2$ while their
Min-SW discrepancy from that Gaussian stays bounded away from zero. The
obstruction is transverse quantization inside projected fibers. This is a
warning about extending an empirical construction to all measures, not an
unqualified statement about every sliced-transport definition.
\source{\bookref{Propositions~11.35--11.36}{230}{thm.11.2.35}}

\post{Why a hard density cap can have no effect on optimal transport}
On a convex Euclidean domain, if both endpoint densities are at most $\kappa$,
the quadratic displacement interpolation also respects this cap. A hard density
constraint therefore leaves the pure Benamou--Brenier optimum unchanged in this
setting. A determinant-concavity argument explains the phenomenon. Obstacles,
nonconvex bottlenecks or additional objectives are needed to produce genuinely
different congestion dynamics.
\illustration{Contrast straight interpolation in a convex domain with transport through a narrow doorway.}
\source{\bookref{Proposition~15.22}{364}{thm.15.2.22}}

\post{Two symmetric central limit theorems, two different Wasserstein rates}
For standardized sums of symmetric Bernoulli variables,
$\Wass_1(\alpha_n,\N(0,1))\sim1/(2\sqrt n)$; for a continuous uniform input,
the leading order is $C/n$. Symmetry removes the skewness contribution in both
cases, but a lattice leaves a first-order sawtooth error in the cumulative
function. This makes the difference between lattice and density asymptotics
visible without changing the limiting Gaussian.
\illustration{Combine the existing rate plot with a magnified cumulative-error profile.}
\source{\bookref{Propositions~10.3--10.4}{188}{thm.10.1.3}}

\post{Why KL belongs to two different divergence families}
Density-ratio divergences and Bregman divergences arise from different
constructions. Under the smoothness and nondegeneracy assumptions of the book,
KL is essentially their common member, up to scaling and irrelevant affine
terms. Comparing second variations forces the Fisher-information form; an
integration then identifies the entropy. The result explains why algorithms for
the two divergence families coincide for KL and generally differ otherwise.
\source{\bookref{Proposition~8.34}{155}{thm.8.8.34}}

\post{Why exact alternating maximization can stop immediately}
Hard $c$-transforms obey a triple-transform identity. Alternating exact best
responses can consequently reach a closed pair without solving the full
Kantorovich dual. The useful lesson is that coordinatewise optimality need not
give global optimality for a nonsmooth concave problem. Entropic smoothing
breaks this locking mechanism and allows repeated updates to keep improving.
\illustration{A small discrete dual with a coordinatewise fixed pair and a better joint direction.}
\source{\bookref{Proposition~5.14}{94}{thm.5.3.14}}

\post{A nonconvex distance kernel produces a geodesically convex MMD}
On the line, take the conditionally positive kernel $k(x,y)=-|x-y|$ with the
energy-distance convention of the book. The self-interaction is affine in
quantile coordinates, whereas the target cross-term is convex. The resulting
squared MMD is displacement convex even though the kernel is not jointly
convex. Ordering on the line, rather than ordinary kernel convexity, is the
decisive mechanism.
\source{\bookref{Proposition~15.18}{360}{thm.15.2.18}}

\post{Sinkhorn slows down when one variable almost determines the other}
The exact local factor for a full Sinkhorn cycle is the square of the maximal
correlation of its optimal coupling. Conditional expectation and the law of total
variance turn an algorithmic rate into a probabilistic statement: nearly
deterministic transport leaves little conditional randomness, hence slow local
convergence. This is a natural sequel to a global Hilbert-metric analysis rather
than a repetition of it.
\illustration{Compare increasingly concentrated couplings and their measured local rates.}
\source{\bookref{Proposition~9.20}{177}{thm.9.5.20}}

\post{Conditional expectation preserves quadratic transport optimality}
If $(X,Y)$ is an optimal coupling for squared Euclidean cost, the map
$x\mapsto\E[Y\mid X=x]$ is optimal from the source to its own output law, even
when the source has atoms. A variance decomposition proves the result. The
qualification about the output law matters: it is generally not the original
target, and the conclusion fails for arbitrary nonoptimal input couplings.
\illustration{Collapse each fan of transport segments onto its conditional barycenter.}
\source{\bookref{Proposition~12.33}{273}{thm.12.6.33}}

\post{A Gaussian target does not make every divergence displacement convex}
Let $\alpha_m=\N(m,1)$ and $\beta=\N(0,1)$. This translation family is a
Wasserstein geodesic, but squared Hellinger divergence satisfies
\[
  H^2(\alpha_m,\beta)=2(1-e^{-m^2/8}),
  \qquad \KL(\alpha_m\mid\beta)=m^2/2.
\]
The Hellinger curvature becomes negative for $|m|>2$, in contrast to KL.
Log-concavity of the target alone cannot transfer KL's displacement-convexity
properties to every $\phi$-divergence.
\illustration{The two scalar functions and their second derivatives.}
\source{\bookref{Remark~15.21}{363}{thm.15.2.21}}

\section*{Further Candidates}
\pdfbookmark[0]{Further Candidates: 13--70}{further}

\post{Composing optimal maps can destroy optimality}
Two symmetric positive-definite linear maps are individually Brenier maps, but
their product need not be symmetric. Noncommuting maps can introduce rotation
through an intermediate measure. This gives an elementary matrix counterexample
to the tempting idea that composing two optimal transport steps is still
optimal. On the line, by contrast, increasing maps remain increasing under
composition.
\source{\bookref{Remark~2.42 and the following example}{30}{thm.2.6.42}}

\post{Polar factorization: sorting a vector field}
Under the appropriate nondegeneracy assumptions, a map factors into a
measure-preserving rearrangement followed by the gradient of a convex function.
The rearrangement changes labels without changing the distribution; the
gradient part carries the distributional change. This is a higher-dimensional
version of sorting, with a useful analogy to the polar decomposition of a
matrix.
\illustration{The original vector field, its relabeling, and its monotone factor.}
\source{\bookref{Proposition~2.28}{22}{thm.2.4.28}}

\post{Triangular rearrangement emerges from an extremely anisotropic metric}
Knothe--Rosenblatt transport is obtained by matching successive conditional
distributions. Under suitable assumptions, it also arises as the limit of
quadratic optimal transport when the coordinate costs have increasingly
separated scales. A singular perturbation of the geometry thus encodes an order
of conditioning.
\illustration{Animate the transport as coordinate anisotropy increases toward a triangular map.}
\source{\bookref{Proposition~2.52}{35}{thm.2.6.52}}

\post{The Gaussian Wasserstein formula is not specifically Gaussian}
The same affine optimal maps and Bures covariance formula apply to elliptical
laws sharing a common radial generator. Gaussianity is sufficient, not
necessary. The shared-generator hypothesis is essential: covariance alone
cannot determine quadratic transport between arbitrary elliptical laws.
\illustration{Compare transport between Gaussian ellipses and between uniform ellipsoids
with the same covariance pairs.}
\source{\bookref{Remark~2.57}{38}{thm.2.7.57}}

\post{Bures distance is a Procrustes problem between matrix square roots}
Writing a covariance as $AA^\top$ leaves an orthogonal ambiguity in its factor.
Minimizing the Frobenius distance between appropriately aligned factors gives
the Bures distance. The same formula links Gaussian transport, shape alignment
and quotient geometry. It is particularly instructive to contrast a fixed
square-root difference with the optimally rotated one.
\source{\bookref{Proposition~2.60}{40}{thm.2.7.60}}

\post{Is a singular covariance nearby or infinitely far away?}
A rank-deficient covariance is at finite Bures distance but lies infinitely far
away in the Gaussian Fisher--Rao geometry. Already in one dimension, standard
deviation gives the Wasserstein coordinate while its logarithm gives the
Fisher--Rao scale coordinate. The same boundary therefore has very different
metric meanings.
\illustration{Compare approaches to the boundary of the two-dimensional PSD cone.}
\source{\bookref{Section~2.7, comparison with Fisher--Rao}{36}{section.2.7}}

\post{Computing quadratic Wasserstein distance without inverting CDFs}
The Bobkov--Ledoux formula expresses one-dimensional squared Wasserstein
distance as a double integral involving positive parts of cumulative-function
differences. A layer-cake argument turns the area associated with squared
displacements into a CDF-only expression. This is an unexpectedly geometric
alternative to the familiar quantile formula.
\illustration{Animate the regions counted by the two integral formulas.}
\source{\bookref{Proposition~2.45}{31}{thm.2.6.45}}

\post{On a tree, optimal transport is bookkeeping across cuts}
For the path metric on a weighted tree, $\Wass_1$ is the sum over edges of the
edge length times the absolute mass imbalance on one side of the cut. Every
unit of imbalance must cross that edge, and the resulting lower bound can be
attained simultaneously. The formula gives a linear-time algorithm and explains
why higher powers of the path distance require a different argument.
\source{\bookref{Proposition~2.50}{33}{thm.2.6.50}}

\post{Removing crossings still leaves exponentially many matchings}
For alternating source and target vertices of a convex polygon, the
noncrossing perfect matchings are counted by Catalan numbers. Eliminating
proper crossings is therefore a powerful geometric simplification but not an
efficient characterization of the optimum: exponentially many candidates can
remain. A short counting argument connects planar transport to classical
combinatorics.
\source{\bookref{Proposition~1.5 and the following Catalan discussion}{7}{thm.1.1.5}}

\post{Concave transport costs reward unequal travel distances}
For concave powers of travel distance, one long trip and one short trip can
cost less than two intermediate trips. The usual intuition behind monotone
quadratic matching is reversed. Simple one-dimensional examples expose nested
matching patterns and show why changing the exponent changes the combinatorics,
not just the numerical values.
\source{\bookref{Section~1.1, concave costs on the line}{1}{section.1.1}}

\post{Independence is almost never optimal transport}
For a continuous cost on compact supports, the independent coupling is optimal
only in the degenerate case where the cost is additive,
$c(x,y)=u(x)+v(y)$ on the product support. In that case every coupling is
optimal. Four-point exchanges prove this rigidity and explain why genuine
transport must create dependence.
\source{\bookref{Proposition~3.29}{56}{thm.3.3.29}}

\post{Optimal transport is the largest convex extension of a pointwise cost}
Prescribe the value $c(x,y)$ on pairs of Dirac masses. Among jointly convex,
lower-semicontinuous extensions with the required Dirac bound, the Kantorovich
value is pointwise largest under the hypotheses stated in the book. Jensen's
inequality applied to a coupling is the key mechanism. This gives an axiomatic
reason for the relaxation, beyond the usual optimization derivation.
\source{\bookref{Theorem~3.35}{59}{thm.3.3.35}}

\post{Transport potentials are always convex in a tropical sense}
Finite pointwise minima of shifted $c$-concave functions remain $c$-concave.
This stability holds even when ordinary convex combinations leave the class.
The proof passes directly through the infimum defining the $c$-transform.
It offers an accessible bridge between optimal transport, min-plus algebra and
architectures built by taking minima of elementary potentials.
\source{\bookref{Proposition~5.17}{95}{thm.5.3.17}}

\post{Convexity of a family of functions can encode curvature of a space}
Whether ordinary averages of $c$-concave potentials remain $c$-concave is tied
to curvature conditions on the cost. Squared Euclidean distance is the simplest
example; spheres and products reveal a less obvious geometric structure.
An expository post can introduce the phenomenon without reproducing the full
regularity theory, using the precise results and assumptions of Kim--McCann and
Figalli--Kim--McCann cited in the book.
\source{\bookref{Section~5.3, paragraph on $c$-concave functions}{91}{section.5.3}}

\post{Bernstein polynomials on the space of probability measures}
Given a functional $f$, form $(B_nf)(\alpha)=\E[f(\hat\alpha_n)]$ from an
empirical measure with $n$ independent samples. This is a polynomial functional
of $\alpha$. Empirical Wasserstein bounds immediately give approximation rates
for Wasserstein--H\"older functionals. The analogy with classical Bernstein
polynomials is exact at the level of sampling and averaging, while the
dimension-dependent rate exposes the difficulty of approximation on measure space.
\source{\bookref{Proposition~4.36 and its preceding definitions}{81}{thm.4.4.36}}

\post{Brenier maps form a closed convex cone}
Fix an absolutely continuous source. Convex combinations of its square-integrable
Brenier maps are still gradients of convex functions and are optimal to their
own pushforward laws. The range of linear optimal-transport coordinates is
closed and convex in the reference $L^2$ space. This rigorously justifies map
interpolation and averaging, without identifying these operations with arbitrary
Wasserstein geodesics or barycenters.
\source{\bookref{Proposition~11.42}{235}{thm.11.4.42}}

\post{Weak changes in a law can control strong changes in its transport map}
Under the source-density and compact convex-domain assumptions stated in the
book, a $\Wass_1^{1/4}$ bound controls the $L^2$ change in the associated
Brenier maps. Stability of the optimal value is much easier than stability of
the optimizer; this theorem bridges the two. A useful post would explain what
the exponent measures and why the reference source must be controlled.
\source{\bookref{Proposition~11.43}{236}{thm.11.4.43}, with the survey and sharper
stability result cited there}

\post{Three Hilbert embeddings, three geometries of distributions}
MMD averages kernel features, sliced Wasserstein records projected quantiles,
and linear OT records reference-to-target maps. All three admit Hilbert-space
descriptions, but they preserve different information and have different image
sets. Consequently, averaging coordinates has a different meaning in each
representation. The comparison makes clear that ``Hilbertian'' is not itself a
complete geometric description.
\source{\bookref{Remark~11.44}{237}{thm.11.4.44}}

\post{Particles and smooth densities have different local sliced geometries}
Along suitable finite atomic curves with fixed weights, the sliced metric
derivative equals the Wasserstein metric derivative divided by $\sqrt d$.
For smooth anisotropic Gaussian deformations, the corresponding inequality can
be strict. Projected ordering explains the atomic equality, while averaging
conditional velocities explains the loss for densities. The contrast warns
against inferring continuum tangent geometry from finite particles alone.
\source{\bookref{Remarks~11.28--11.29}{225}{thm.11.2.28}}

\post{Averaging projected distributions need not produce a distribution}
Independently averaging quantiles in each projection direction need not produce
a Radon-consistent family. Reconstructing the result by a pseudoinverse may even
leave the cone of nonnegative densities. The surprise is that individually valid
one-dimensional distributions need not fit together into any multidimensional
law. This connects sliced barycenters to tomography and compatibility constraints.
\source{\bookref{Proposition~12.15 and its surrounding Radon discussion}{255}{thm.12.1.15}}

\post{Optimal quantization uses averages of quantiles, not midpoint quantiles}
For $n$ equally weighted atoms on the line, the optimal location of atom $j$ is
$n\int_{(j-1)/n}^{j/n}C_\alpha^{-1}(r)\,dr$, rather than the midpoint quantile
in general. The proof is an orthogonal projection of the quantile function onto
piecewise constant functions. This gives a precise distinction between optimal
deterministic quantization and a common sampling heuristic.
\source{\bookref{Proposition~6.14}{110}{thm.6.4.14}}

\post{Carefully placed particles beat random particles by a full power}
For regular one-dimensional laws satisfying the relevant quantile assumptions,
optimal equal-weight quantization has squared Wasserstein error of order
$n^{-2}$, while the expected squared error of independent particles is of order
$n^{-1}$. Quantile approximation gives the first rate; a Brownian-bridge limit
gives the second. The statement should retain the regularity and integrability
conditions, rather than suggesting a universal rate for all one-dimensional laws.
\source{\bookref{Propositions~6.14--6.15}{110}{thm.6.4.14}}

\post{A discrete barycenter can be much sparser than its candidate support}
For $S$ discrete input laws with support sizes $n_s$, a barycenter exists with
at most $\sum_s n_s-S+1$ atoms. This can be dramatically smaller than the
product-sized candidate set suggested by all tuples of input points. The proof
is a rank count for a multimarginal linear program, connecting transport
sparsity to basic feasible solutions.
\source{\bookref{Corollary~12.22}{260}{thm.12.2.22}}

\post{Multimarginal transport meets graphical models}
When a multimarginal cost factors along a graph, the Gibbs tensor also factors.
Message passing can then compute Sinkhorn contractions without constructing
the full tensor. The decisive exponential dependence is governed by treewidth,
not necessarily by the number of marginals. Trees, cycles and dense graphs make
the contrast concrete and connect transport to graphical models and tensor networks.
\source{\bookref{Section~12.2 and Example~12.24}{257}{section.12.2}}

\post{Sinkhorn divergences forget the reference measure}
Changing a compatible product reference in entropic OT adds marginal KL terms
to the value. These terms cancel exactly when the two self-costs are subtracted
to form the Sinkhorn divergence. The result explains why apparently different
entropy conventions can yield the same debiased discrepancy, provided the
absolute-continuity and finiteness assumptions are satisfied.
\source{\bookref{Remark~8.38}{159}{thm.8.9.38}}

\post{An entropic transport problem requiring no Sinkhorn iterations}
Choose a Gibbs transition from the source, normalize it appropriately, and take
its output law as the target. The resulting coupling already solves the
entropic transport problem. For quadratic cost and a Lebesgue reference, this
includes Gaussian noising. The example links a closed-form Schr\"odinger scaling
to the noising paths used in diffusion models.
\source{\bookref{Proposition~8.14}{138}{thm.8.4.14}}

\post{Large-temperature Sinkhorn projects onto interaction terms}
The first correction to the independent coupling is governed by the cost after
subtracting its source average and target average and restoring its global
average. This double centering is the interaction term of an ANOVA
decomposition. A large-temperature expansion thus has a simple Hilbert-space
interpretation: row-only and column-only information cannot change the coupling.
\source{\bookref{Proposition~8.16}{140}{thm.8.4.16}}

\post{Fisher information hides in the small-noise transport bias}
For sufficiently regular quadratic transport, the second-order entropic
correction, after the appropriate endpoint and normalization terms are removed,
involves Fisher information integrated along the displacement interpolation.
The bias therefore detects the roughness of intermediate densities, not just
the endpoint displacement. A careful post should state the normalization before
displaying the expansion: it determines which terms appear at each order.
\source{\bookref{Proposition~8.17}{141}{thm.8.4.17}}

\post{Why half the temperature can give a better barycenter}
In a doubly regularized barycenter, an outer entropy with strength
$\tau=\epsilon/2$ cancels the first-order candidate-entropy contribution of the
inner entropic transport costs, under the small-temperature expansion's
hypotheses. A seemingly arbitrary parameter choice thus has an asymptotic
explanation. Contrast sharpening, entropy cancellation and uniqueness, which
are related but distinct issues.
\source{\bookref{Definition~12.13 and the subsequent discussion}{252}{thm.12.1.13}}

\post{Eliminating a Sinkhorn potential improves conditioning computably}
The semi-dual Hessian at the optimum is $(I-T^*T)/\epsilon$, where $T$ is
conditional expectation under the optimal coupling. Variable projection turns
the full Hessian into a Schur complement. Because its diagonal blocks are
identities after normalization, the spectral conditioning bounds can be compared
explicitly. Distinguish a sharp comparison of available bounds from an
unconditional equality between actual condition-number improvement factors.
\source{\bookref{Proposition~9.24}{181}{thm.9.5.24}}

\post{Sinkhorn acceleration is classical SOR in disguise}
Over-relaxing each of Sinkhorn's two block updates gives a nonlinear analogue
of successive over-relaxation. Linearization explains the optimal parameter and
the square-root improvement in a small local spectral gap. This provides a
clean reason for acceleration without introducing an unrelated momentum model.
The local nature of the optimal-rate statement should remain explicit.
\illustration{Reuse the marginal-error curves together with predicted and measured gains.}
\source{\bookref{Proposition~9.22}{179}{thm.9.5.22}}

\post{What happens when the Sinkhorn temperature becomes complex?}
Near a positive real temperature, gauge-fixed Sinkhorn scalings have a local
holomorphic continuation under the hypotheses of the continuation theorem.
The implicit-function mechanism is familiar, but its application to couplings
produces striking oscillatory images. The continuation is local: it is not a
license to substitute any complex temperature into a real formula or to assume
global convergence of complex Sinkhorn iterations.
\source{\bookref{Theorem~8.43}{161}{thm.8.10.43}, with the
Carlier--Pegon--Tamanini attribution given there}

\post{Allowing mass variation can make Sinkhorn convergence simpler}
KL penalties on the marginals introduce damping factors into unbalanced
Sinkhorn updates. They break the balanced problem's additive gauge and yield a
direct contraction in the potential sup norm in the stated bounded setting.
The result contrasts with the projective geometry needed for balanced scaling:
relaxing the constraints can simplify the convergence mechanism.
\source{\bookref{Proposition~11.12}{217}{thm.11.1.12}}

\post{Changing the regularizer changes whether transport has exact zeros}
Entropy produces an exponential primal--dual law and hence dense couplings
under finite-cost positivity assumptions. A quadratic regularizer instead
produces thresholding and can create exact zeros. Generalized soft
$c$-transforms connect this distinction to one-dimensional monotone equations
and simplex projections. Sparsity here is an analytic consequence of the
regularizer, not merely a numerical threshold.
\source{\bookref{Section~8.8 and Proposition~8.30}{153}{section.8.8}}

\post{Unbalanced transport separates shape from a scalar mass problem}
With equal KL marginal penalties of strength $\tau$, write the input masses as
$A$ and $B$ and optimize first over a normalized coupling shape. If its optimal
shape cost is $E$, the transported mass is
$M^*=\sqrt{AB}\exp(-E/(2\tau))$ in the book's normalization. Geometry enters a
one-dimensional mass minimization as an exponential discount. The geometric
mean of the input masses is the unpenalized mass scale.
\source{\bookref{Proposition~11.3}{210}{thm.11.1.3}}

\post{The Hellinger interpolation of two Gaussians has three components}
Hellinger interpolation is linear in square-root densities. Squaring a
combination of two Gaussian square roots gives the two endpoint densities and
a Gaussian geometric-mean cross-term: an unnormalized three-Gaussian mixture.
This is a transparent counterexample to Gaussian closure for pure
Fisher--Rao/Hellinger motion. Its relation to WFR should be phrased as a limiting
obstruction, not as an explicit formula for every WFR geodesic.
\source{\bookref{Remark~11.18}{219}{thm.11.1.18}}

\post{Partial transport trades a price against transported mass}
The optimal cost at prescribed transported mass is a convex value function.
Adding a linear price for unmatched mass selects a supporting slope of this
curve. Trimming arguments connect this formulation to TV-penalized unbalanced
transport under the assumptions of the equivalence result. The scalar value
curve makes clear when several transported masses can share the same price.
\illustration{Animate the active transport region as the price changes.}
\source{\bookref{Proposition~11.21}{220}{thm.11.1.21}}

\post{Forgetting mixture labels is a Wasserstein contraction}
A probability law over probability laws can be collapsed to its mixture.
Couple the outer laws and then glue suitable inner couplings: the resulting
coupling proves that collapse is non-expansive for Wasserstein-over-Wasserstein.
The proof identifies precisely what is forgotten when hierarchical labels are
discarded. It is also a compact example of gluing as the engine of a metric
comparison.
\source{\bookref{Proposition~13.10}{286}{thm.13.2.10}}

\post{Compare shapes through distributions of distance distributions}
Associate to each point the law of its distance to a random point. A shape then
becomes a distribution over one-dimensional distributions. M\'emoli's profile
bound compares these objects and lower-bounds intrinsic GW. Conditional
couplings explain the inequality, while the construction connects directly to
Wasserstein-over-Wasserstein.
\illustration{Reuse the cat--bunny figure with matched points and colored distance profiles.}
\source{\bookref{Proposition~13.24}{292}{thm.13.3.24}}

\post{A nonnegative distortion can define a concave coupling objective}
Squared GW distortion is pointwise nonnegative, yet its quadratic form can be
concave on zero-marginal directions when the dissimilarities have the required
conditional definiteness. A tensor-feature calculation turns the second
variation into a negative squared norm. This conditional concavity explains
extreme-point optimizers and permutation solutions in appropriate finite
uniform settings, without making the optimization convex or globally easy.
\source{\bookref{Proposition~13.27}{295}{thm.13.3.27}}

\post{Euclidean GW can be understood as learning a cost matrix}
For inner-product or squared-distance dissimilarities, completing a square
reduces the learned-cost representation of GW to a finite-dimensional matrix
parameter with a quadratic penalty. The remaining OT problem uses a cost
depending on that matrix. This interpretation separates cost learning from
coupling computation and reveals why these Euclidean models are unusually
tractable, although the outer optimization is still nonconvex.
\source{\bookref{Proposition~13.32}{298}{thm.13.3.32}}

\post{Duplicating a variable can exactly relax a nonconvex problem}
Replacing the diagonal quadratic GW objective by a two-coupling bilinear
objective gives a lower bound and a problem convex in either block separately.
Konno's polarization argument shows tightness in the conditionally concave
regime. With positive entropy, strict convexity additionally constrains
minimizing block pairs. The distinction between exactness of the relaxation
and global convergence of alternating minimization is an essential part of the story.
\source{\bookref{Proposition~13.39}{303}{thm.13.3.39}}

\post{Weak transport measures distance to a mean-preserving spread}
Quadratic barycentric weak transport equals the minimum of
$\Wass_2(\alpha,\eta)^2$ over $\eta\preceq_{\mathrm{cx}}\beta$. It vanishes
exactly when a martingale coupling exists, subject to the standard moment
assumptions. Conditional Jensen gives one direction; the convex-order
characterization supplies the other. The result connects a convex transport
relaxation to mean-preserving spreads and mathematical finance.
\source{\bookref{Proposition~12.44}{278}{thm.12.6.44}}

\post{The right topology does not guarantee the triangle inequality}
A quantum-inspired classical transport cost compares square-root coupling
densities with their transposes. It is convex in the coupling, nonnegative and
separating; Wasserstein comparisons give the desired topology in the stated
setting. Nevertheless its square-root version fails the triangle inequality
already on three points. The example separates topological usefulness from
metric structure particularly cleanly.
\source{\bookref{Propositions~13.53 and 13.56, Example~13.57}{310}{thm.13.4.53}}

\post{Replacing scalar scalings by matrices changes Sinkhorn's meaning}
Classical row and column scalings commute in ways that matrix-valued updates do
not. In quantum OT, operator scaling is generally not the same as exact
Bregman projection for the entropically regularized transport objective. The
diagonal case hides this distinction because it reduces to classical Sinkhorn.
A two-by-two noncommuting example could explain why copying a scalar formula
does not automatically copy its variational interpretation.
\source{\bookref{Section~13.4, Gurvits scaling and quantum Sinkhorn}{308}{section.13.4}}

\post{Particles can move while their probability law stands still}
A rotating cloud can have a stationary radial distribution. Its particle
velocity is therefore not the intrinsic velocity of the measure curve. The
minimum-kinetic-energy representative removes invisible circulation and is a
gradient, characterized by a weighted Poisson problem. This is a tangible
explanation of the quotient hidden in Wasserstein tangent geometry.
\source{\bookref{Proposition~14.3}{322}{thm.14.1.3}}

\post{The same perspective transform convexifies transport and divergences}
The momentum action $|m|^2/\rho$ and the divergence integrand
$q\phi(p/q)$ are instances of the same perspective construction. Its convexity
explains both the Benamou--Brenier momentum formulation and joint convexity of
$\phi$-divergences. Concave mobilities and unbalanced actions extend the same
mechanism. A single geometric proof can replace several apparently unrelated
convexity calculations.
\source{\bookref{Proposition~14.13 and its cross-reference to divergence convexity}{332}{thm.14.4.13}}

\post{Why random-path optimization reduces to endpoint optimization}
The KL chain rule splits path-space entropy into endpoint entropy and the
average conditional-bridge entropy. At fixed endpoints, the optimum keeps the
reference bridges unchanged. The Schr\"odinger problem consequently reduces to
an entropic coupling problem, followed by a bridge lift. This gives a clean
reason why a finite-dimensional-looking Sinkhorn problem can describe a law on
an infinite-dimensional path space.
\source{\bookref{Proposition~14.11}{329}{thm.14.3.11}}

\post{The logarithmic mean is forced by linear Markov diffusion}
To make entropy descent produce a linear jump flux, the edge mobility must
satisfy $\theta(a,b)(\log a-\log b)=a-b$. This identifies the logarithmic mean,
including its continuous extension on the diagonal. A seemingly specialized
mean is therefore selected by a chain rule. The same identity connects discrete
Markov chains and continuum jump kernels.
\illustration{An edge with unequal densities, its entropy gradient, and the resulting linear flux.}
\source{\bookref{Proposition~15.58 and the following Markov-chain discussion}{388}{thm.15.6.58}}

\post{Fractional heat is entropy descent in a different transport geometry}
Replace local velocities by antisymmetric jump fluxes and use a reversible
power-law jump kernel. Entropy descent then produces a fractional diffusion
operator rather than the ordinary Laplacian. The energy can remain familiar
while the transport geometry changes. An accessible post can derive the flux
identity before discussing the fractional heat equation.
\source{\bookref{Section~15.6}{388}{section.15.6}}

\post{One heat equation, two gradient-flow structures}
The heat equation is both $L^2$ descent of Dirichlet energy and Wasserstein
descent of entropy. A gradient flow is therefore specified by an energy and a
geometry together, not by the PDE alone. Deriving the same equation twice gives
a short and useful introduction to metric-dependent steepest descent.
\illustration{Compare what a unit tangent displacement means in the two geometries.}
\source{\bookref{Remark~15.14}{354}{thm.15.1.14}}

\post{A spectral gap is a linearized entropy inequality}
Perturb a minimizer of a Wasserstein--PL inequality and retain the quadratic
terms. For relative entropy, logarithmic Sobolev linearizes to Poincar\'e, which
is a spectral gap for the target-weighted Laplacian. The mechanism connects
nonlinear entropy decay, local quadratic forms and linear operator theory.
This can serve as a sequel to a PL post rather than another abstract
convergence theorem.
\source{\bookref{Proposition~15.44}{376}{thm.15.3.44} and
\bookref{Proposition~15.48}{377}{thm.15.3.48}}

\post{A convex functional of measures can have bad particle equilibria}
The functional $f(\alpha)=\int V\,d\alpha$ is affine in the measure even if
$V$ has nonglobal local minima. Particles can become trapped at those minima.
Convexity under mixing measures is not convexity under moving their support.
This elementary example explains why affine measure convexity alone is not a
guarantee of global convergence for particle training.
\source{\bookref{Proposition~15.51 and its surrounding discussion}{381}{thm.15.4.51}}

\post{Mean--covariance closure does not require Gaussianity}
If $f(\alpha)=g(m_\alpha,\Sigma_\alpha)$, its Wasserstein descent velocity is
affine in space for every input law. Mean and covariance therefore obey an
autonomous system even for non-Gaussian initialization. Gaussianity is only what
turns these two statistics into a complete description of the law; other
initial shapes undergo an affine deformation. This cleanly separates moment
closure from distributional closure.
\source{\bookref{Equation~(16.43), distribution-free moment closure}{425}{equation.16.5.43}}

\post{Smooth random paths and Brownian paths can share all time marginals}
A static Gaussian noise variable reused along an interpolant can produce the
same one-time laws as an Ornstein--Uhlenbeck noising process, despite completely
different temporal dependence and quadratic variation. Matching marginals does
not identify a path-space law. The example is an effective antidote to
confusing a prescribed distributional path with a unique stochastic process.
\source{\bookref{Proposition~16.9}{405}{thm.16.1.9}}

\post{Gaussian flow matching detects noncommuting covariances}
For the independent Gaussian interpolation and scalar schedules studied in the
book, the resulting terminal transport map is Brenier precisely under the
stated covariance commutation condition. Noncommutativity introduces a
geometric obstruction that changing the scalar schedule cannot generally
remove. This is a sharper test of transport optimality than simply checking
that the endpoint marginals are correct.
\source{\bookref{Proposition~16.7}{404}{thm.16.1.7}}

\post{Gaussian blurring mean shift eventually makes one cluster}
For blurring mean shift with a positive Gaussian kernel and compactly supported
initial data, the invariant convex hull gives a uniform overlap bound, and a
Dobrushin contraction yields exponential consensus. Several visible modes can
therefore be long transients rather than permanent clusters. This contrasts
with fixed-data mean shift and with bounded-confidence kernels, for which the
same global consensus conclusion need not hold.
\source{\bookref{Remark~16.19}{416}{thm.16.4.19}}

\post{Unbiased samples do not give unbiased transport gradients}
An empirical measure inserted inside a nonlinear transport optimization does
not generally give an unbiased estimate of the population gradient. This is
different from the usual stochastic-gradient story for an average of losses.
A multimodal example can show how minibatches alter both the coupling and the
induced update, rather than adding only centered noise to a fixed gradient.
\source{\bookref{Section~16.2, one-step models and minibatch transport}{407}{section.16.2}}

\clearpage
\section*{Longer Posts and Follow-Ups}
\pdfbookmark[0]{Longer Posts and Follow-Ups: 71--75}{followups}

\post{The basic GW iteration is a difference-of-convex algorithm}
In the conditionally concave regime, the coupling objective has a global affine
upper model at each iterate. Minimizing that frozen OT model is a classical
DC/majorization step; Toland duality supplies a complementary cost-side view.
The connection clarifies descent and fixed-point conditions without confusing
stationarity with global optimality. It also distinguishes this interpretation
from alternating minimization of an arbitrary biconvex relaxation.
\source{\bookref{Remark~13.31}{297}{thm.13.3.31}}

\post{Optimal dual potentials determine statistical fluctuations}
On finite supports, with unique normalized dual potentials, the two-sample
limiting variance is
\[
  \operatorname{Var}_{\alpha}(f^*)+
  \operatorname{Var}_{\beta}(g^*)
\]
for equal sample sizes and the book's normalization. Dual potentials act as
influence functions. Nonuniqueness can instead produce non-Gaussian limits.
The post should explain why fixed support sizes matter, rather than treating
this CLT as a dimension-free theorem for arbitrary continuous transport.
\source{\bookref{Proposition~10.18}{200}{thm.10.3.18}}

\post{Gaussian attention, mean shift and denoising share one identity}
For Gaussian kernel averaging with variance $\epsilon$, the normalized local
mean satisfies
\[
  m_\epsilon[\alpha](x)-x
  =\epsilon\nabla\log\rho_\epsilon[\alpha](x).
\]
Differentiating the Gaussian convolution is enough to prove it. The same
identity interprets a weighted average as a score step, connects attention to
mean shift, and gives the denoising interpretation associated with Tweedie's
formula. The kernel convention must be fixed to get the factor $\epsilon$ right.
\source{\bookref{Equation~(16.41), attention and mean shift}{416}{equation.16.4.41}}

\post{Training a generator does not define an intrinsic distributional flow}
Two parameterizations representing the same law can induce different particle
and measure velocities under parameter gradient descent. The Eulerian velocity
is a conditional expectation of parameter-induced particle motion. This formula
pinpoints where the parameterization enters and explains why Euclidean training
of a generator is not automatically Wasserstein steepest descent of its output
law.
\source{\bookref{Proposition~16.12}{408}{thm.16.2.12}}

\post{An irrelevant additive constant becomes a growth rate}
For probability measures, first variations are defined only modulo additive
constants, because admissible mass variations have total mass zero. In an
unbalanced transport geometry, the same constant affects reaction and therefore
total mass. The passage from transport to transport--reaction turns a harmless
gauge into a genuine dynamical parameter. A one-line energy shift gives a
particularly clear illustration.
\source{\bookref{Section~15.7}{390}{section.15.7}}

\end{document}
